\documentclass[12pt,a4paper]{article}
\usepackage[utf8]{inputenc}
\usepackage[T1]{fontenc}
\usepackage{geometry}
\usepackage{amsmath}
\usepackage{booktabs}
\usepackage{graphicx}
\usepackage{float}
\usepackage{xcolor}
\usepackage{listings}
\geometry{margin=2.5cm}
\graphicspath{{./}}
\renewcommand{\arraystretch}{1.2}
\renewcommand{\figurename}{Figura}
\renewcommand{\tablename}{Tabela}
\definecolor{fundocodigo}{RGB}{247,247,247}
\lstset{basicstyle=\ttfamily\footnotesize,backgroundcolor=\color{fundocodigo},
  frame=single,breaklines=true,columns=fullflexible,showstringspaces=false,
  literate={á}{{\'a}}1 {é}{{\'e}}1 {í}{{\'i}}1 {ó}{{\'o}}1 {ú}{{\'u}}1
           {ã}{{\~a}}1 {õ}{{\~o}}1 {ç}{{\c c}}1}

\title{Questão 2 -- Caracterização e Análise Descritiva do Conjunto de Dados\\
  \large Dados de Compartilhamento de Bicicletas}
\author{}
\date{}

\begin{document}
\maketitle

\section{Introdução}

A Questão 2 analisa o conjunto específico do grupo construído na Questão 1.
São 300 observações diárias consecutivas, identificadas pelos índices 7 a 306,
entre 7 de janeiro e 2 de novembro de 2011. A variável
\texttt{total\_user} registra o número diário de usuários e corresponde à soma
de \texttt{casual} e \texttt{registered}. A análise começa pela estrutura e
completude da base, passa pelas medidas de tendência central e pelos quartis,
examina a distribuição em gráficos e termina com a classificação dos dias de
baixa utilização.

Todos os resultados numéricos a seguir se referem à amostra de 300 dias, e não
à base original de 731 observações. Nas seções correspondentes, são indicados
o procedimento, os resultados e a interpretação de cada etapa.

\section{Classificação e identificação das variáveis (etapa 2.1)}

\subsection{Procedimento}

Inicialmente, foram identificadas as colunas da planilha e verificada a
função de cada uma. A coluna \texttt{...1} reproduz o índice
\texttt{instant}; por isso, é apenas uma coluna auxiliar e não foi contada
como uma variável adicional. Entre as oito variáveis de interesse,
\texttt{instant} é um identificador numérico discreto, \texttt{dteday} é uma
data, \texttt{season} e \texttt{weathersit} são categóricas, e
\texttt{temp}, \texttt{casual}, \texttt{registered} e
\texttt{total\_user} são quantitativas. As três últimas contagens são
discretas; a temperatura é contínua.

\begin{table}[H]
\centering\footnotesize
\renewcommand{\arraystretch}{1.06}
\caption{Natureza e interpretação das variáveis da amostra.}
\begin{tabular}{p{2.3cm}p{3.6cm}p{8.6cm}}
\toprule
Variável & Natureza & Interpretação e unidade \\
\midrule
\texttt{instant} & Índice discreto & Identificador sequencial, de 7 a 306; sem unidade física. \\
\texttt{dteday} & Temporal & Data da observação; resolução de um dia. \\
\texttt{season} & Categórica & Estação do ano: inverno, primavera, verão ou outono. \\
\texttt{weathersit} & Categórica & Condição meteorológica registrada no dia. \\
\texttt{temp} & Contínua & Temperatura na escala numérica da planilha, de 2,4 a 34,8. \\
\texttt{casual} & Discreta & Contagem diária de usuários casuais, em usuários. \\
\texttt{registered} & Discreta & Contagem diária de usuários registrados, em usuários. \\
\texttt{total\_user} & Discreta derivada & Contagem diária total, em usuários. \\
\bottomrule
\end{tabular}
\end{table}

No segundo passo, contaram-se as categorias observadas em
\texttt{season} e \texttt{weathersit}. Foram encontrados 73 dias de inverno,
92 de primavera, 94 de verão e 41 de outono. Para as condições
meteorológicas, a amostra contém 184 dias de céu limpo, 106 dias nublados e
10 dias de chuva fraca. As frequências de cada variável somam 300.

Por fim, verificou-se a presença de valores ausentes em cada coluna: não há
\texttt{NA} em nenhuma das 300 linhas. Conferiu-se também a relação
\begin{equation}
\texttt{total\_user}_i=\texttt{casual}_i+\texttt{registered}_i,
\end{equation}
válida para todas as observações. Essa verificação é necessária porque o
total de usuários será usado nas etapas seguintes.

\subsection{Conclusão da etapa}

A amostra tem oito variáveis analíticas, uma coluna auxiliar de indexação e
nenhum valor ausente. As estações e as condições meteorológicas estão
representadas pelas categorias indicadas acima. A ausência de valores
faltantes dispensa imputação nesta análise, mas, isoladamente, não comprova
ausência de outros problemas de qualidade de dados.

\section{Medidas de tendência central (etapa 2.2)}

\subsection{Seleção das variáveis e cálculo da média}

Foram selecionadas \texttt{temp}, \texttt{casual}, \texttt{registered} e
\texttt{total\_user}, pois representam medidas quantitativas interpretáveis.
Não se calcula a média de \texttt{instant}, que apenas identifica registros,
nem dos códigos das categorias. Para cada variável $x$, a média das
$n=300$ observações é
\begin{equation}
\overline{x}=\frac{1}{300}\sum_{i=1}^{300}x_i.
\end{equation}
Por exemplo, a soma de \texttt{total\_user} é 1.053.964; logo,
\begin{equation}
\overline{\texttt{total\_user}}=
\frac{1.053.964}{300}=3.513{,}2133\simeq 3.513{,}21.
\end{equation}

\subsection{Ordenação, mediana e moda}

Em seguida, ordenaram-se separadamente os valores de cada variável em ordem
crescente. Como 300 é par, a mediana é a média dos valores nas posições 150 e
151 da sequência ordenada:
\begin{equation}
\operatorname{Md}(x)=\frac{x_{(150)}+x_{(151)}}{2}.
\end{equation}
Para \texttt{total\_user}, esses valores são 3.982 e 4.010, resultando em
mediana de 3.996. Para a moda, contaram-se as ocorrências de cada valor e
identificou-se a maior frequência. A moda, diferentemente da mediana, não
precisa ser única.

\begin{table}[H]
\centering\small
\caption{Medidas de tendência central nas 300 observações.}
\begin{tabular}{lrrp{5.5cm}}
\toprule
Variável & Média & Mediana & Moda \\
\midrule
\texttt{temp} & 21,15 & 22,2 & 26,0 (6 ocorrências) \\
\texttt{casual} & 743,77 & 673 & 639 (3 ocorrências) \\
\texttt{registered} & 2.769,45 & 3.093 & 16 valores empatados, cada um com 2 ocorrências \\
\texttt{total\_user} & 3.513,21 & 3.996 & 9 valores empatados, cada um com 2 ocorrências \\
\bottomrule
\end{tabular}
\end{table}

\subsection{Conclusão da etapa}

Os usuários registrados representam a maior parcela do uso: sua média é
2.769,45, contra 743,77 para usuários casuais. A mediana do total (3.996)
supera sua média (3.513,21), mostrando que o valor médio é sensível aos dias
de menor procura. Para \texttt{registered} e \texttt{total\_user}, há várias
modas empatadas com frequência máxima de apenas duas observações; portanto,
uma moda isolada não descreve bem o centro dessas distribuições. A relação
entre média e mediana sugere uma distribuição não simétrica, mas a forma
completa deve ser examinada nos gráficos da etapa 2.4.

\section{Quartis e identificação de valores atípicos (etapa 2.3)}

\subsection{Ordenação e cálculo dos quartis}

A variável escolhida foi \texttt{total\_user}, pois reúne os usuários
casuais e registrados e mede a demanda total em cada dia. Antes de calcular
os quartis, seus 300 valores foram colocados em ordem crescente. Foi usado o
método \texttt{type = 7} de \texttt{quantile()} no R. Nesse método, a posição
teórica do quantil de probabilidade $p$ é
\begin{equation}
h=1+(n-1)p.
\end{equation}
Para $p=0{,}25$, $h=75{,}75$; para $p=0{,}75$, $h=225{,}25$.
Valores em posições fracionárias são obtidos por interpolação entre as
observações ordenadas vizinhas. Os resultados são
\begin{equation}
Q_1=2.105{,}5,\qquad Q_2=3.996,\qquad Q_3=4.681.
\end{equation}

O procedimento pode ser reproduzido no R por:
\begin{lstlisting}[language=R]
x <- sort(dados$total_user)
quartis <- quantile(x, probs = c(0.25, 0.50, 0.75), type = 7)
print(quartis)
\end{lstlisting}

\subsection{Intervalo interquartil e limites}

A dispersão dos 50\% centrais dos valores é medida pelo intervalo
interquartil:
\begin{equation}
IQR=Q_3-Q_1=4.681-2.105{,}5=2.575{,}5.
\end{equation}
Aplicando a regra de $1{,}5\,IQR$, obtêm-se
\begin{align}
L_{\mathrm{inf}}&=Q_1-1{,}5\,IQR=-1.757{,}75,\\
L_{\mathrm{sup}}&=Q_3+1{,}5\,IQR=8.544{,}25.
\end{align}
Um valor inferior ao primeiro limite ou superior ao segundo seria marcado
como potencialmente atípico. Como o mínimo é 431 e o máximo é 6.043, nenhuma
das 300 observações ultrapassa esses limites.

\subsection{Conclusão da etapa}

O primeiro quartil delimita a região dos menores valores e servirá de corte
para a variável \texttt{low\_usage}. O IQR mostra uma diferença de 2.575,5
usuários entre o primeiro e o terceiro quartis. Não há outliers segundo a
regra de $1{,}5\,IQR$ na distribuição global. Isso não significa que todos
os dias tenham demandas semelhantes nem permite concluir, por si só, que a
demanda seja previsível.

\section{Histograma e boxplot (etapa 2.4)}

\subsection{Construção e leitura do histograma}

\begin{figure}[H]
\centering
\includegraphics[width=.72\textwidth]{histogramaq2.png}
\caption{Histograma de densidade de \texttt{total\_user} fornecido para a etapa 2.4.}
\label{fig:hist}
\end{figure}

Com os 300 totais diários, construiu-se um histograma de densidade: o eixo
horizontal representa usuários por dia, e o vertical representa densidade,
não quantidade absoluta de dias. Cada barra reúne valores em uma faixa de
utilização. Na Figura~\ref{fig:hist}, observam-se concentrações aproximadas
nas faixas de 1.500 a 2.000 e de 4.500 a 5.000 usuários. A aparência de
dois agrupamentos é descritiva e pode variar com a largura das classes; sua
causa não pode ser determinada apenas pelo histograma.

\subsection{Construção e leitura do boxplot}

O boxplot resume os mesmos dados com uma caixa entre $Q_1$ e $Q_3$, uma
linha na mediana e hastes que alcançam os extremos não classificados como
outliers pela regra de $1{,}5\,IQR$. A Figura~\ref{fig:box} mostra mediana
mais próxima de $Q_3$ do que de $Q_1$, bem como maior extensão da parte
inferior da caixa. Não se veem pontos fora das hastes, em concordância com
os limites calculados na etapa 2.3.

\begin{figure}[H]
\centering
\includegraphics[width=.74\textwidth]{boxplotq2.png}
\caption{Boxplot de \texttt{total\_user} fornecido para a etapa 2.4.}
\label{fig:box}
\end{figure}

\subsection{Conclusão da etapa}

Os gráficos se complementam. O histograma permite observar onde os valores
se concentram e sugere dois grupos de intensidade de uso. O boxplot evidencia
os quartis, a amplitude e a ausência de observações atípicas sob a regra
adotada. Nenhum dos dois gráficos, sem comparações adicionais, demonstra que
as diferenças sejam causadas por estação do ano, clima ou tipo de dia.

\section{Criação da variável de baixa utilização (etapa 2.5)}
\vspace{0.35em}

\subsection{Definição e implementação}

Tomou-se o primeiro quartil da distribuição global de
\texttt{total\_user} como limiar. Para cada dia $i$, definiu-se
\begin{equation}
\texttt{low\_usage}_i=
\begin{cases}
1,&\texttt{total\_user}_i<2.105{,}5,\\
0,&\texttt{total\_user}_i\geq 2.105{,}5.
\end{cases}
\end{equation}
Como \texttt{total\_user} é uma contagem inteira, o primeiro caso equivale
a ter no máximo 2.105 usuários. Os dados originais não são alterados; apenas
se acrescenta a nova coluna binária. O comando usado para reproduzir esse
passo em R é:
\begin{lstlisting}[language=R]
Q1 <- quantile(dados$total_user, probs = 0.25, type = 7)
dados$low_usage <- ifelse(dados$total_user < Q1, 1L, 0L)
table(dados$low_usage)
mean(dados$low_usage)
\end{lstlisting}

\subsection{Contagem e proporção}

Após aplicar o corte às 300 linhas, 75 dias receberam valor 1 e 225 dias
receberam valor 0. A proporção classificada como baixa utilização é
\begin{equation}
\widehat{p}=\frac{75}{300}=0{,}25=25\%.
\end{equation}

\begin{table}[H]
\centering
\caption{Distribuição da variável \texttt{low\_usage}.}
\begin{tabular}{lrr}
\toprule
Classificação & Dias & Proporção \\
\midrule
Baixa utilização (1) & 75 & 25\% \\
Demais dias (0) & 225 & 75\% \\
\midrule
Total & 300 & 100\% \\
\bottomrule
\end{tabular}
\end{table}

\subsection{Conclusão da etapa}

Na amostra estudada, baixa utilização significa pertencer ao grupo dos
menores totais diários segundo o corte do primeiro quartil: menos de
2.105,5 usuários. O resultado de 25\% decorre do critério relativo à
distribuição do próprio grupo. Esses dias não são automaticamente outliers,
nem se pode transportar o mesmo limiar para outro período sem reavaliar a
comparabilidade dos dados.

Em síntese, a média de \texttt{total\_user} foi 3.513,21, inferior à
mediana de 3.996, e o corte em $Q_1$ identificou 75 dias de baixa demanda.
O histograma sugere concentrações distintas de uso, sem que os gráficos ou
os quartis permitam atribuir causas a esse padrão. Investigar possíveis
associações com estação do ano, clima ou calendário exige análises
específicas posteriores.

\end{document}
